\documentclass[xcolor=dvipsnames]{beamer}
%\usepackage[OT4]{fontenc}
%\usepackage[utf8]{inputenc}
%\usepackage[polish]{babel}
\usepackage{fancyvrb}
\usepackage{thumbpdf}
\usepackage{relsize}
\usepackage{amsmath}

%\usetheme[left]{Marburg}
%\usetheme{Hannower}
\useinnertheme[shadow]{rounded}
\useoutertheme[right,width=2cm,hideothersubsections]{sidebar}

\definecolor{ZooplusGreen}{RGB}{129,197,46}
\definecolor{ZooplusBlue}{RGB}{4,6,76}
\setbeamercolor{structure}{fg=ZooplusGreen}
\setbeamercolor{palette primary}{fg=ZooplusBlue,bg=ZooplusGreen!70}
\setbeamercolor{palette secondary}{fg=ZooplusBlue,bg=ZooplusGreen!80}
\setbeamercolor{palette tertiary}{fg=ZooplusBlue,bg=ZooplusGreen!90}
\setbeamercolor{palette quaternary}{fg=ZooplusBlue,bg=ZooplusGreen}
\setbeamercolor{titlelike}{parent=palette quaternary}
\setbeamercolor{block title}{fg=ZooplusBlue,bg=ZooplusGreen}
\setbeamercolor{block title alerted}{use=alerted text,fg=ZooplusBlue,bg=alerted text.fg!75!bg}
\setbeamercolor{block title example}{use=example text,fg=ZooplusBlue,bg=example text.fg!75!bg}
\setbeamercolor{block body}{parent=normal text,use=block title,bg=block title.bg!25!bg}
\setbeamercolor{block body alerted}{parent=normal text,use=block title alerted,bg=block title alerted.bg!25!bg}
\setbeamercolor{block body example}{parent=normal text,use=block title example,bg=block title example.bg!25!bg}
\setbeamercolor{sidebar}{bg=ZooplusGreen!70}
\setbeamercolor{palette sidebar primary}{fg=ZooplusBlue}
\setbeamercolor{palette sidebar secondary}{fg=ZooplusBlue!75}
\setbeamercolor{palette sidebar tertiary}{fg=ZooplusBlue!75}
\setbeamercolor{palette sidebar quaternary}{fg=ZooplusBlue}
\setbeamercolor*{separation line}{}
\setbeamercolor*{fine separation line}{}
\logo{\includegraphics[scale=0.25]{../../zooplus_logo.png}}

\usefonttheme{default}
\setbeamercovered{transparent}
\title{Akka -- actors system for Scala \& Java}
\author{Jacek~Bilski}
\date{\today}
\subject{Akka -- actors system for Scala \& Java}

\setbeamertemplate{navigation symbols}
{
	\usebeamercolor[fg]{navigation symbols dimmed}
	{
		\insertframenumber\,/\,\inserttotalframenumber
	}
}

\begin{document}

\begin{frame}
\titlepage
\end{frame}

\begin{frame}
\frametitle{Agenda}
\tableofcontents[pausesections]
\end{frame}

\section{Akka}

\begin{frame}
\frametitle{Akka}
\begin{itemize}
\item \href{http://akka.io/}{http://akka.io/}
\item Implements actors system for Scala \& Java
\item Part of Scala right now (well, almost)
\end{itemize}
\end{frame}


\section{Actors system}

\begin{frame}
\frametitle{Actors system}
\begin{itemize}
\item Hierarchy of actors
\item Manages actors, creates threads, routes messages, /etc.
\end{itemize}
\end{frame}

\begin{frame}
\frametitle{Actors}
\begin{itemize}
\item Very lightweight concurrent entities
\item Process messages asynchronously using an event-driven receive loop
\item ''Little'' SOA
\item No need to worry about threads, locks, /etc.
\end{itemize}
\end{frame}

\section{Problem}

\begin{frame}
\frametitle{Problem: calculating value of $\pi$}
Formula:
%\begin{equation}
\[\sum_{n=0}^\infty \frac{(-1)^n}{2n+1} = 1 - \frac{1}{3} + \frac{1}{5} - \frac{1}{7} + \ldots = \frac{\pi}{4}\]
%\end{equation}
\end{frame}

\section{Solution}

\begin{frame}
\frametitle{Divide and conquer}
\begin{itemize}
\item Problem is easily parallelizable,\ldots
\item \ldots{}so we split calculations in chunks and calculate in different actors.
\end{itemize}
\end{frame}

\frame[containsverbatim]{
\frametitle{Worker}
\begin{Verbatim}[obeytabs=true,fontsize=\relsize{-2},tabsize=2]
class Worker extends Actor {

def calculatePiFor(start: Int, nrOfElements: Int): BigDecimal = {
	var acc = 0.0
	for (i <- start until (start + nrOfElements))
		acc += 4.0 * (1 - (i % 2) * 2) / (2 * i + 1)
	acc
}

def receive = {
	case Work(start, nrOfElements) =>
		// perform the work
		sender ! Result(calculatePiFor(start, nrOfElements))
}
}
\end{Verbatim}
}

\frame[containsverbatim]{
\frametitle{Master}
\begin{Verbatim}[obeytabs=true,fontsize=\relscale{0.7},tabsize=2]
class Master extends Actor {

var pi: BigDecimal = 0
var nrOfResults: Int = _

val workerRouter = context.actorOf(
	Props[Worker].withRouter(RoundRobinRouter(nrOfWorkers)),
		name = "workerRouter")

def receive = {
	case Calculate =>
		for (i <- 0 until nrOfMessages)
			workerRouter ! Work(i * nrOfElements, nrOfElements)
	case Result(value) =>
		pi += value
		nrOfResults += 1
		if (nrOfResults == nrOfMessages) {
			// Send the result to the listener
			listener ! PiApproximation(pi)
			// Stops this actor and all its supervised children
			context.stop(self)
		}
}}
\end{Verbatim}
}

\frame[containsverbatim]{
\frametitle{Listener}
\begin{Verbatim}[obeytabs=true,fontsize=\relscale{0.8},tabsize=2]
class Listener extends Actor {
def receive = {
	case PiApproximation(pi, duration) =>
		println("\n\tPi approximation: \t\t%s".format(pi))
		println("\n\tCalculation time: \t%s".format(duration))
		context.system.shutdown()
}}
\end{Verbatim}
}

\frame[containsverbatim]{
\frametitle{App}
\begin{Verbatim}[obeytabs=true,fontsize=\relscale{0.7},tabsize=2]
calculate(workersNo = 16, elementsNo = 10000, messagesNo = 100000)

def calculate(workersNo: Int, elementsNo: Int, messagesNo: Int) {
	// Create an Akka system
	val system = ActorSystem("PiSystem")

	// create the result listener, which will print the result and
	// shutdown the system
	val listener = system.actorOf(Props[Listener], name = "listener")

	// create the master
	val master = system.actorOf(Props(new Master(
		workersNo, messagesNo, elementsNo, listener)),
		name = "master")

	// start the calculation
	master ! Calculate
}
\end{Verbatim}
}

\begin{frame}
\frametitle{Results}
$\pi = 3.141592652589784499499528724894 \gets$ no wonder :)
Execution times:\\
1 worker -- 13275 milliseconds\\
4 workers -- 4323 milliseconds\\
8 workers -- 3734 milliseconds\\
16 workers -- 3612 milliseconds
\end{frame}

\section{Questions?}

\begin{frame}
\frametitle{Questions?}
\begin{center}
\Huge{?}
\end{center}
\end{frame}


\end{document}
